% Reporte del proyecto final DevStudio
% Compilar desde esta carpeta: /usr/bin/latexbuild reporte.tex
% Preámbulo compartido por el reporte, los manuales y el contrato.
% Se compila con XeLaTeX mediante /usr/bin/latexbuild.
\documentclass[11pt,a4paper]{article}

\usepackage{fontspec}
\setmonofont{JetBrainsMono Nerd Font Mono}[Scale=0.82]
\usepackage[spanish,es-nodecimaldot,es-noshorthands]{babel}

\usepackage[top=2.4cm,bottom=2.4cm,left=2.4cm,right=2.4cm]{geometry}
\usepackage{fancyhdr}
\usepackage{titlesec}
\usepackage{xcolor}
\usepackage{booktabs}
\usepackage{tabularx}
\usepackage{array}
\usepackage{enumitem}
\usepackage{float}
\usepackage{graphicx}
\usepackage{tikz}
\usetikzlibrary{positioning,arrows.meta}
\usepackage{listings}
\usepackage{xurl}

\makeatletter
\renewcommand*\l@section[2]{%
  \ifnum \c@tocdepth >\z@
    \addpenalty\@secpenalty
    \addvspace{1.0em \@plus\p@}%
    \setlength\@tempdima{2.6em}%
    \begingroup
      \parindent \z@ \rightskip \@pnumwidth
      \parfillskip -\@pnumwidth
      \leavevmode \bfseries
      \advance\leftskip\@tempdima
      \hskip -\leftskip
      #1\nobreak\hfil \nobreak\hb@xt@\@pnumwidth{\hss #2}\par
    \endgroup
  \fi}
\makeatother
\usepackage[hidelinks]{hyperref}

\definecolor{tnBlue}{HTML}{0EA5E9}
\definecolor{tnDark}{HTML}{0F172A}
\definecolor{codeGray}{HTML}{F4F6F8}
\definecolor{codeGreen}{HTML}{2E7D32}
\definecolor{codeRed}{HTML}{C62828}

\newcolumntype{Y}{>{\raggedright\arraybackslash}X}

\lstset{
  basicstyle=\ttfamily\footnotesize,
  breaklines=true,
  breakatwhitespace=false,
  frame=tb,
  framerule=0.4pt,
  rulecolor=\color{tnDark!40},
  framesep=5pt,
  keywordstyle=\color{tnBlue}\bfseries,
  commentstyle=\color{codeGreen}\itshape,
  stringstyle=\color{codeRed},
  showstringspaces=false,
  tabsize=2,
  columns=fullflexible,
  keepspaces=true,
  extendedchars=true,
  captionpos=t,
  xleftmargin=0pt,
  xrightmargin=0pt,
  aboveskip=0.9em,
  belowskip=0.9em,
  literate=
    {á}{{\'a}}1 {é}{{\'e}}1 {í}{{\'i}}1 {ó}{{\'o}}1 {ú}{{\'u}}1
    {Á}{{\'A}}1 {É}{{\'E}}1 {Í}{{\'I}}1 {Ó}{{\'O}}1 {Ú}{{\'U}}1
    {ñ}{{\~n}}1 {Ñ}{{\~N}}1 {ü}{{\"u}}1 {Ü}{{\"U}}1
    {¿}{{?`}}1 {¡}{{!`}}1
    {·}{{\textperiodcentered}}1
}

\pagestyle{fancy}
\fancyhf{}
\fancyhead[L]{\small Alfredo Ramírez Candia}
\fancyhead[R]{\small Diseño de Aplicaciones Web}
\fancyfoot[C]{\thepage}
\renewcommand{\headrulewidth}{0.4pt}

\titleformat{\section}
  {\Large\bfseries\color{tnDark}}{\thesection.}{0.5em}{}
\titleformat{name=\section,numberless}
  {\Large\bfseries\color{tnDark}}{}{0em}{}
\titleformat{\subsection}
  {\large\bfseries\color{tnBlue}}{\thesubsection}{0.5em}{}
\titlespacing*{\section}{0pt}{1.4em}{0.6em}
\titlespacing*{\subsection}{0pt}{1.05em}{0.35em}

\addto\captionsspanish{%
  \renewcommand{\tablename}{Tabla}%
  \renewcommand{\figurename}{Figura}%
  \renewcommand{\lstlistingname}{Listado}%
}

\newcommand{\datosPdf}[1]{%
  \hypersetup{%
    pdftitle={#1},%
    pdfauthor={Alfredo Ramírez Candia},%
    pdflang={es}%
  }%
}

\newcommand{\portada}[3]{%
  \begin{titlepage}
    \centering
    \vspace*{0.8cm}
    {\large\bfseries BENEMÉRITA UNIVERSIDAD AUTÓNOMA DE PUEBLA\par}
    \vspace{0.28cm}
    {\large FACULTAD DE CIENCIAS DE LA COMPUTACIÓN\par}
    \vspace{0.28cm}
    {\large LICENCIATURA/INGENIERÍA EN CIENCIAS DE LA COMPUTACIÓN\par}
    \vspace{1.35cm}
    \rule{\linewidth}{0.5pt}\par
    \vspace{0.45cm}
    {\Huge\bfseries #1\par}
    \vspace{0.38cm}
    {\Large #2\par}
    \vspace{0.32cm}
    {\large Asignatura: Diseño de Aplicaciones Web\par}
    \vspace{0.45cm}
    \rule{\linewidth}{0.5pt}\par
    \vspace{1.5cm}
    \begin{tabular}{rl}
      \textbf{Docente:}   & Miguel Ángel Vargas Lomelí \\[0.55cm]
      \textbf{Alumno:}    & Alfredo Ramírez Candia \\[0.55cm]
      \textbf{Matrícula:} & 202375947 \\[0.55cm]
      \textbf{Fecha:}     & 22 de noviembre de 2026 \\
    \end{tabular}
    \vfill
    {\small #3\par}
  \end{titlepage}
  \setcounter{page}{2}
}

\setcounter{tocdepth}{2}
\setcounter{secnumdepth}{2}
\setlist[enumerate]{leftmargin=1.5em,itemsep=0.45em,topsep=0.4em}
\setlist[itemize]{leftmargin=1.4em,itemsep=0.28em,topsep=0.35em}
\emergencystretch=2em
\urlstyle{tt}
\renewcommand{\arraystretch}{1.18}

\datosPdf{Proyecto final. DevStudio}

\newcommand{\captura}[4][width=\linewidth]{%
  \begin{figure}[H]
    \centering
    \includegraphics[#1]{figuras/#2}%
    \caption{#3}%
    \label{#4}%
  \end{figure}%
}

\begin{document}
\portada{Proyecto final}{DevStudio, sitio y portal para negocios de Puebla}{Proyecto final de la asignatura}

\tableofcontents
\clearpage

\section{Presentación}

DevStudio es la aplicación web del proyecto final. Es el sitio de un estudio
que trabaja con negocios de Puebla: el visitante entiende el problema, el
resultado y el siguiente paso, y la tecnología queda en segundo plano. El
cliente registrado consulta sus proyectos, el avance y el calendario. El
administrador atiende solicitudes, mensajes y el contenido.

El trabajo es individual. El código vive en \path{proyecto-final/}: el frente
estático en \path{appWeb/} y la API en \path{backEnd/}. Esta carpeta
\path{reporte/} es la memoria del proyecto. Los manuales y el contrato están
en carpetas hermanas y no sustituyen a la documentación Markdown previa.

\subsection{Alcance}

El sitio público presenta soluciones, casos, la forma de trabajo, contacto y
la solicitud de diagnóstico. El registro crea un cliente. El portal muestra
solo los proyectos y las solicitudes de esa cuenta, permite pedir o cancelar
reuniones y cambiar la foto de perfil. El panel cubre usuarios, clientes,
servicios, tecnologías, proyectos, solicitudes y mensajes. Confirmar una
reunión o escribir el avance queda en la API de administración: no hay todavía
una pantalla de desarrolladores para esas dos acciones.

Quedan fuera los pagos, la aplicación móvil nativa, el correo transaccional y
el hosting de producción. Esas exclusiones se repiten en el contrato.

Las figuras~\ref{fig:inicio} a~\ref{fig:diagnostico} muestran ese sitio en
marcha. El menú público no cambia al iniciar sesión: la cuenta queda a la
derecha.

\captura{01-inicio.png}{Inicio en tema de día, con la sesión de la cliente de demostración.}{fig:inicio}

\captura{02-inicio-noche.png}{El mismo inicio en tema de noche.}{fig:noche}

\captura[width=0.38\linewidth]{03-inicio-movil.png}{Inicio a 375 píxeles. El menú pasa a una segunda línea y el llamado a la acción sigue a la vista.}{fig:movil}

\captura{04-servicios.png}{Listado de soluciones.}{fig:servicios}

\captura{05-caso.png}{Caso publicado: portal de clientes de Comercial Andina.}{fig:caso}

\captura{06-diagnostico.png}{Página de diagnóstico, con las referencias de inversión.}{fig:diagnostico}

\subsection{Tecnologías}

El manual de PHP (PHP Documentation Group, s.~f.), la documentación de
SQL Server (Microsoft, s.~f.) y Docker Compose (Docker, Inc., s.~f.)
sostienen la pila. Este PDF se compiló con XeTeX (TeX Users Group, s.~f.)
y con latexmk (Collins, 2025).

\begin{table}[H]
  \centering
  \caption{Piezas del sistema y el papel de cada una.}
  \label{tab:tecnologias}
  \small
  \begin{tabularx}{\textwidth}{@{}l Y@{}}
    \toprule
    Pieza & Papel \\
    \midrule
    HTML y CSS & Páginas estáticas, mobile first, variables de día y noche \\
    JavaScript & Sesión en el encabezado, formularios y calendario \\
    PHP 8.3 & API REST, sesión, CSRF y reglas de negocio \\
    SQL Server 2022 & Persistencia, restricciones e historial \\
    Nginx & Sirve el frente y reenvía \path{/api} \\
    Docker Compose & Levanta web, API y base en un solo comando \\
    \bottomrule
  \end{tabularx}
\end{table}

\section{Modelación}

\subsection{Requisitos}

Hay tres actores. El visitante consulta el sitio y envía contacto o
diagnóstico sin cuenta. El cliente, rol \texttt{cliente}, ve y edita su
perfil, consulta sus proyectos y solicitudes, agenda reuniones y no puede
escribir el porcentaje de avance. El administrador, rol \texttt{admin},
gestiona catálogos, bandejas y, por la API, el progreso y la respuesta a las
reuniones.

Los requisitos funcionales que el sistema ya cumple son estos:

\begin{enumerate}
  \item Presentar la empresa, las soluciones y los casos publicados.
  \item Recibir contacto y diagnóstico con validación y límite de envíos.
  \item Registrar, iniciar y cerrar sesión, con roles distintos.
  \item Impedir que un cliente lea o cambie datos de otra cuenta.
  \item Mostrar el avance de 0 a 100 y el historial de notas.
  \item Consultar, agendar y cancelar reuniones del propio cliente.
  \item Confirmar o rechazar reuniones y registrar avances desde la API de administración.
  \item Cambiar la foto de perfil y conservarla al navegar.
\end{enumerate}

Los requisitos no funcionales que condicionan el diseño son el mismo origen
para la cookie de sesión, el diseño usable en 375, 768 y 1280 píxeles, las
consultas parametrizadas y la reconstrucción de la base desde los scripts.

\subsection{Mapa de sitio}

El frente no usa URLs limpias. Nginx entrega el HTML tal como está en disco.
La tabla~\ref{tab:mapa} resume las páginas reales.

\begin{table}[H]
  \centering
  \caption{Mapa de sitio implementado.}
  \label{tab:mapa}
  \small
  \begin{tabularx}{\textwidth}{@{}Y l@{}}
    \toprule
    Página & Quién entra \\
    \midrule
    \path{/index.html} & Visitante \\
    \path{/pages/servicios.html} y detalle por \texttt{slug} & Visitante \\
    \path{/pages/portafolio.html} y detalle de caso & Visitante \\
    \path{/pages/nosotros.html}, contacto y cotización & Visitante \\
    \path{/pages/login.html} y \path{/pages/registro.html} & Visitante \\
    \path{/pages/portal/index.html} & Cliente \\
    Proyectos, detalle, solicitudes y detalle & Cliente \\
    \path{/pages/portal/calendario.html} & Cliente \\
    \path{/pages/portal/perfil.html} & Cliente \\
    \path{/pages/admin/} dashboard, usuarios, clientes & Admin \\
    Servicios, proyectos, tecnologías, solicitudes, mensajes & Admin \\
    \bottomrule
  \end{tabularx}
\end{table}

\subsection{Guiones}

Cuatro recorridos bastan para explicar el uso. El visitante entra al inicio,
abre una solución y envía el diagnóstico en \path{/pages/cotizacion.html}. Si
faltan datos, el formulario conserva lo escrito. La solicitud queda en estado
\texttt{nueva}.

Quien crea cuenta captura nombre, correo, empresa y una contraseña de al menos
ocho caracteres, con mayúscula, minúscula y número. El sistema lo deja en el
login. Al entrar, el cliente va al portal y el administrador al panel. Las
figuras~\ref{fig:login} y~\ref{fig:registro} son esas dos pantallas, todavía
sin sesión.

\captura{07-login.png}{Inicio de sesión. Desde aquí se crea la cuenta de cliente.}{fig:login}

\captura{08-registro.png}{Alta de cliente. El rol de administrador no se elige en este formulario.}{fig:registro}

Laura, la cliente de demostración, abre Mis proyectos y ve la barra de avance.
En Calendario elige un hueco de entre 15 minutos y 4 horas. La reunión nace
como \texttt{solicitada}. Puede cancelarla si aún no empieza. No puede
escribir los puntos ni la respuesta: eso lo hace el desarrollador por la API.
Las figuras~\ref{fig:resumen} a~\ref{fig:perfil} recorren su portal.

\captura{09-resumen.png}{Resumen de Laura Méndez. La foto abre el perfil.}{fig:resumen}

\captura{10-proyectos.png}{Proyectos de la cuenta, con el avance en 15\,\% y 45\,\%.}{fig:proyectos}

\captura{11-proyecto.png}{Detalle del portal de Comercial Andina y la nota del 18 de noviembre.}{fig:proyecto}

\captura{12-calendario.png}{Calendario de noviembre de 2026 y el formulario para pedir una reunión.}{fig:calendario}

\captura{13-perfil.png}{Perfil de la cliente, con la foto ampliada.}{fig:perfil}

El administrador revisa la bandeja de solicitudes y mensajes. Puede convertir
una solicitud en proyecto. El porcentaje de avance se actualiza con
\path{PATCH /api/v1/admin/proyectos/{id}/progreso}. Un cliente que llame esa
ruta recibe 403. Las figuras~\ref{fig:admin} a~\ref{fig:mensajes} son el panel
con el que entra \texttt{admin@devstudio.local}.

\captura{14-admin.png}{Resumen del panel. El monograma del administrador es AD.}{fig:admin}

\captura{15-solicitudes.png}{Bandeja de solicitudes, con estado, plan y monto.}{fig:solicitudes}

\captura{16-mensajes.png}{Bandeja de mensajes de contacto.}{fig:mensajes}

\subsection{Diagrama de navegación}

La figura~\ref{fig:navegacion} separa el sitio público de las dos zonas con
sesión. El encabezado no cambia de liga al entrar: el acceso a la cuenta está
en el monograma, a la derecha del menú.

\begin{figure}[H]
  \centering
  \begin{tikzpicture}[
      nodo/.style={
        draw=tnDark!75, line width=0.6pt, rounded corners=2pt,
        fill=codeGray, align=center, inner sep=4pt, font=\footnotesize,
        text width=3.1cm
      },
      flecha/.style={-{Stealth[length=2mm]}, line width=0.6pt, draw=tnDark}
    ]
    \node[nodo] (inicio) {Inicio};
    \node[nodo, right=8mm of inicio] (sol) {Soluciones y casos};
    \node[nodo, right=8mm of sol] (cot) {Contacto y diagnóstico};
    \node[nodo, below=12mm of inicio] (login) {Login y registro};
    \node[nodo, below=12mm of sol] (portal) {Portal del cliente};
    \node[nodo, below=12mm of cot] (admin) {Panel de administración};
    \draw[flecha] (inicio) -- (sol);
    \draw[flecha] (sol) -- (cot);
    \draw[flecha] (inicio) -- (login);
    \draw[flecha] (login) -- (portal);
    \draw[flecha] (login) -- (admin);
  \end{tikzpicture}
  \caption{Del sitio público a las dos áreas con sesión.}
  \label{fig:navegacion}
\end{figure}

Dentro del portal el orden es Resumen, Proyectos, Solicitudes, Calendario y
Perfil. La foto del resumen abre Perfil. Abajo de la actividad reciente está
Cerrar sesión.

\subsection{Estructura de archivos}

\begin{table}[H]
  \centering
  \caption{Carpetas que importan para instalar y entregar.}
  \label{tab:carpetas}
  \small
  \begin{tabularx}{\textwidth}{@{}l Y@{}}
    \toprule
    Ruta & Contenido \\
    \midrule
    \path{appWeb/} & HTML, CSS y JavaScript del sitio, portal y panel \\
    \path{backEnd/public/} & Único punto de entrada de la API \\
    \path{backEnd/app/} & Controladores, modelos y sesión \\
    \path{backEnd/database/} & \path{schema.sql}, \path{seed.sql} y migraciones \\
    \path{docker/} & Imagen PHP y configuración de Nginx \\
    \path{compose.yaml} & Web, API y SQL Server \\
    \path{reporte/} y hermanos & Esta documentación en LaTeX \\
    \bottomrule
  \end{tabularx}
\end{table}

\section{Arquitectura y entorno}

\subsection{Diagrama}

El navegador solo habla con Nginx. Así la cookie de sesión viaja en el mismo
origen y no hace falta CORS. La figura~\ref{fig:arquitectura} muestra el
recorrido.

\begin{figure}[H]
  \centering
  \begin{tikzpicture}[
      caja/.style={
        draw=tnDark!75, line width=0.7pt, rounded corners=3pt,
        fill=codeGray, align=center, inner sep=7pt, text width=8.6cm,
        font=\small
      },
      flecha/.style={-{Stealth[length=2.3mm]}, line width=0.75pt, draw=tnDark}
    ]
    \node[caja] (nav) {Navegador};
    \node[caja, below=7mm of nav] (ngx)
      {Nginx, puerto 8080\\[1pt]
       {\footnotesize HTML de \texttt{appWeb} y proxy de \texttt{/api}}};
    \node[caja, below=7mm of ngx] (php)
      {PHP 8.3 y Apache\\[1pt]
       {\footnotesize sesión, CSRF y \texttt{pdo\_sqlsrv}}};
    \node[caja, below=7mm of php] (sql)
      {SQL Server 2022\\[1pt]
       {\footnotesize base \texttt{DevStudioDB}, puerto 1433}};
    \draw[flecha] (nav) -- (ngx);
    \draw[flecha] (ngx) -- (php);
    \draw[flecha] (php) -- (sql);
  \end{tikzpicture}
  \caption{Mismo origen: el navegador no llama a la API por otro host.}
  \label{fig:arquitectura}
\end{figure}

\subsection{Puertos}

\begin{table}[H]
  \centering
  \caption{Servicios de \texttt{compose.yaml}.}
  \label{tab:puertos}
  \begin{tabular}{lll}
    \toprule
    Servicio & Contenedor & Puerto publicado \\
    \midrule
    web & \texttt{pf-web} & 8080 \\
    api & \texttt{pf-api} & solo red interna \\
    mssql & \texttt{pf-mssql} & 1433 \\
    \bottomrule
  \end{tabular}
\end{table}

La zona horaria de negocio es \texttt{America/Mexico\_City}. Las reuniones se
guardan en UTC y el calendario las muestra en esa zona.

\section{Desarrollo}

\subsection{HTML y CSS, mobile first}

Las páginas comparten encabezado, pie y variables de color. El tema de día o
noche se aplica en una línea de script antes de las hojas de estilo, así no
hay un destello claro al abrir de noche. La figura~\ref{fig:noche} es ese
tema y la figura~\ref{fig:movil} es el mismo inicio a 375 píxeles. La
cuadrícula del calendario y las tarjetas del portal se probaron también en
768 y 1280 píxeles. El menú público sigue siendo Soluciones, Casos, Cómo
trabajamos y Contacto. La cuenta y el cambio de tema viven aparte, a la
derecha.

\subsection{JavaScript, usuarios y restricciones}

El módulo de interfaz lee la sesión en \path{/api/v1/auth/me} y muestra el
monograma o el enlace de entrar. Las iniciales y una miniatura de la foto se
guardan en \texttt{localStorage} para pintarse antes del primer render: al
cambiar de página la foto no vuelve a pedirse a la red. Las páginas del portal
exigen rol \texttt{cliente}; las del panel, rol \texttt{admin}. Esa guarda es
de interfaz. La API repite la comprobación.

\subsection{PHP}

\path{backEnd/public/index.php} compara el método y la ruta. Los identificadores
de la URL tienen que ser dígitos. Cada escritura, salvo el alta pública de
contacto y diagnóstico, exige el encabezado \texttt{X-CSRF-Token}. El cliente
de una ruta \path{/mi/} sale de la sesión, nunca del cuerpo. Un administrador
en esas rutas recibe 403, y un cliente en \path{/admin/} también.

El avance acepta un entero de 0 a 100 y una nota de 2 a 500 caracteres. Guarda
el porcentaje en el proyecto y añade una fila al historial. Una reunión nueva
queda \texttt{solicitada}. Solo se cancela si está solicitada o confirmada y
todavía no empieza. No puede empalmarse con otra reunión activa del mismo
cliente.

\subsection{Base de datos}

El esquema es idempotente. Las tablas que sostienen el negocio son
\texttt{Roles}, \texttt{Usuarios}, \texttt{Clientes}, \texttt{Servicios},
\texttt{Tecnologias}, \texttt{Proyectos}, \texttt{ProyectoTecnologia},
\texttt{ProyectoAvances}, \texttt{Reuniones}, \texttt{Planes},
\texttt{SolicitudesCotizacion} y \texttt{MensajesContacto}.

\texttt{Proyectos.progreso} es un \texttt{TINYINT} de 0 a 100.
\texttt{ProyectoAvances} guarda porcentaje, nota y fecha.
\texttt{Reuniones} exige inicio anterior al fin y una duración de 15 a 240
minutos. El estado solo puede ser \texttt{solicitada}, \texttt{confirmada},
\texttt{rechazada} o \texttt{cancelada}. \texttt{Usuarios.foto} guarda la
extensión del archivo, que vive fuera del directorio público.

\section{Resultados}

\subsection{Matriz de cumplimiento}

\begin{table}[H]
  \centering
  \caption{Requisitos de la materia y dónde quedan cubiertos.}
  \label{tab:matriz}
  \small
  \begin{tabularx}{\textwidth}{@{}Y Y@{}}
    \toprule
    Requisito & Dónde se cumple \\
    \midrule
    Requisitos, mapa, guiones, navegación y archivos & Secciones 2.1 a 2.5 de este reporte \\
    HTML, CSS mobile first y JavaScript con login & \path{appWeb/} y la sección 4 \\
    PHP y diseño de la base & \path{backEnd/} y la sección 4 \\
    Manual de usuario & \path{manual-usuario/build/manual-usuario.pdf} \\
    Manual técnico & \path{manual-tecnico/build/manual-tecnico.pdf} \\
    Contrato de prestación, alcances y costo & \path{contrato/build/contrato.pdf} \\
    Código en ZIP y PDF en LaTeX & \path{scripts/empaquetar.sh} \\
    \bottomrule
  \end{tabularx}
\end{table}

\subsection{Cuentas de demostración}

El seed crea dos cuentas para recorrer el sistema. No se documenta aquí la
contraseña del servidor de base de datos.

\begin{table}[H]
  \centering
  \caption{Cuentas locales después de \texttt{seed.sql}.}
  \label{tab:cuentas}
  \begin{tabular}{lll}
    \toprule
    Rol & Correo & Contraseña \\
    \midrule
    admin & \texttt{admin@devstudio.local} & \texttt{Admin\#2026} \\
    cliente & \texttt{cliente@devstudio.local} & \texttt{Cliente\#2026} \\
    \bottomrule
  \end{tabular}
\end{table}

La cliente de demostración es Laura Méndez, de Comercial Andina. Su calendario
no abre vacío: hay un avance y dos reuniones de ejemplo.

\section{Cómo reproducir}

Hace falta Docker y los puertos 8080 y 1433 libres. Desde
\path{proyecto-final/}:

\begin{lstlisting}[caption={Levantar el entorno.}]
cp .env.example .env
docker compose up -d --build
\end{lstlisting}

Hay que definir \texttt{DB\_PASSWORD} en \path{.env} antes de arrancar. Luego
se cargan el esquema y los datos:

\begin{lstlisting}[caption={Cargar esquema y semilla.}]
docker compose exec -T mssql /bin/bash -c \
  '/opt/mssql-tools18/bin/sqlcmd -S localhost -U sa -P "$MSSQL_SA_PASSWORD" -C -i /dev/stdin' \
  < backEnd/database/schema.sql
docker compose exec -T mssql /bin/bash -c \
  '/opt/mssql-tools18/bin/sqlcmd -S localhost -U sa -P "$MSSQL_SA_PASSWORD" -C -i /dev/stdin' \
  < backEnd/database/seed.sql
\end{lstlisting}

El sitio queda en \url{http://localhost:8080}. El estado de la API se consulta
en \path{/api/v1/health}. Los cuatro PDF de esta entrega se generan con
\path{/usr/bin/latexbuild} dentro de cada carpeta, o de una vez con
\path{scripts/empaquetar.sh}.

\section{Conclusiones}

El proyecto cubre el recorrido que pide la materia: modelar, construir con
HTML, CSS y JavaScript, proteger el acceso en PHP y dejar la base diseñada y
reproducible. El cliente ya puede ver su avance y usar el calendario. Lo que
sigue fuera de estas páginas, a propósito, es la pantalla en la que el
desarrollador confirmaría la reunión y escribiría el avance: la API ya lo
permite y el contrato lo deja dentro del alcance del sistema, no de una
entrega posterior de interfaz.

\section{Referencias}

\begin{list}{}{%
  \setlength{\leftmargin}{1.27cm}%
  \setlength{\itemindent}{-1.27cm}%
  \setlength{\itemsep}{0.7em}%
  \setlength{\parsep}{0pt}%
  \setlength{\topsep}{0.4em}%
}
  \item Collins, J. (2025). \emph{latexmk} (Versión 4.87) [Software]. \url{https://ctan.org/pkg/latexmk}

  \item Docker, Inc. (s.~f.). \emph{Docker Compose}. \url{https://docs.docker.com/compose/}

  \item Microsoft. (s.~f.). \emph{Documentación técnica de SQL Server}. \url{https://learn.microsoft.com/es-es/sql/sql-server/sql-server-technical-documentation}

  \item PHP Documentation Group. (s.~f.). \emph{Manual de PHP}. \url{https://www.php.net/manual/es/}

  \item TeX Users Group. (s.~f.). \emph{XeTeX}. \url{https://tug.org/xetex/}
\end{list}

\end{document}
